% \begin{Details}

\subsection{Motivation}


\subsubsection{Applications}
Shear links are used in the Richmond-San Rafel Bridge \citep{seible2000long}.

The response of bearings \cite{montalto2024buckling,montalto2024effective}.

Several finite element analysis products implement distinct models for special cases like thin-walled or open sections.
\subsection{Objectives}

While the trace formalism clarifies the classical shear correction debate, the 
primary objective of this work is not the determination of shear correction 
factors \emph{per se}. Rather, our goal is to develop a framework for obtaining 
\emph{warping modes} suitable for inelastic beam analysis.

In linear elastic analysis, the distinction between different traces is 
largely academic: all traces that exactly represent the Saint-Venant class 
yield equivalent global stiffness matrices, differing only in the physical 
interpretation assigned to boundary conditions. The situation changes 
fundamentally when material nonlinearity is introduced. 
Inelastic constitutive 
evaluation requires knowledge of the actual strain distribution at material 
points throughout the cross-section—not merely the integrated stress 
resultants. This distribution depends on the \emph{warping shapes} employed 
in the kinematic description, and the classical Saint-Venant warping functions 
do not, in general, furnish the correct shapes for a given trace.

The connection to shear correction arises because addressing this requirement 
forces us to confront the same fundamental ambiguity that underlies the shear 
correction debate. Different traces imply different associations between 
beam strains and pointwise deformations, and hence different effective 
warping shapes. The trace matrix $\TraceDeDpRoot$ encodes precisely this 
association: it transforms the canonical Ie\c{s}an parameters of the 
Saint-Venant solution into the strain measures of a particular beam theory, 
and its inverse determines the warping shapes that reconstruct the continuum 
strain field from those beam strains.

A secondary objective is computational efficiency. Existing implementations 
of inelastic beam elements with warping \citep{saritas2009inelastic,%
lecorvec2012nonlinear,dire2018mixed} introduce auxiliary kinematic fields 
that must be interpolated along the element length. These additional degrees 
of freedom either inflate the global system of equations—incurring a cost 
that scales as the cube of the DOF count for direct linear solvers—or must 
be condensed out at the element level during each Newton iteration. We seek 
a formulation in which warping effects are entirely absorbed into the 
section-level constitutive relation, so that standard six-DOF beam elements 
can be employed without modification.

\subsection{Scope}

The developments in this work are restricted to the following setting:
\begin{itemize}
\item Prismatic, initially straight beams with arbitrary cross-section geometry.
\item Homogeneous, isotropic material response at the continuum level 
      (though the section-level formulation accommodates inelasticity 
      through fiber discretization).
\item The Saint-Venant class of loadings: axial force, biaxial bending moments, 
      torque, and biaxial shear forces, all applied as end resultants.
\item Quasi-static response; dynamic and frequency-matching considerations 
      are outside our scope.
\item Small to moderate material strain.
\end{itemize}

We do not treat thin-walled sections as a special case requiring separate 
theoretical development. The formulation applies to arbitrary solid or 
thin-walled cross-sections, with the classical thin-walled results 
(Vlasov torsion, etc.) emerging as special cases.

\subsection{Outline}

The remainder of this work is organized as follows:

\Cref{sec:genwarp} develops the governing equations for a geometrically exact 
beam enriched with an arbitrary number of warping modes. 
Beam theories are classified into four types according to how warping amplitudes are related to the primary beam strains, and show that the constitutive relations for all types are connected by linear transformations. 
The Type~3 formulation, in which warping amplitudes are algebraically constrained to the beam strains, 
is identified as the target for standard six-DOF finite elements.

\Cref{sec:gensec} develops the section-level constitutive theory that connects 
the one-dimensional beam model to the three-dimensional continuum. 
We introduce material points in the cross-section through an extended displacement field and derive strain measures consistent with the beam kinematics.

\Cref{sec:timo} introduces the trace formalism and characterizes the family 
of traces satisfying our admissibility criteria. We derive explicit expressions 
for the trace matrices corresponding to classical beam theories and introduce 
the energy-conjugate trace, which produces symmetric section stiffness and arises naturally from variational considerations.

Finally, we present numerical examples demonstrating exact recovery of 
Saint~Venant solutions through the proposed formulation and illustrating 
the behavior under inelastic material response.
% \end{Details}

\subsection{Overview}

The draft first lays out the nonlinear governing equations of equilibrium for general lines in space in \cref{sec:genwarp}. 
These equations are determined solely from the geometric nature of a line embedded in space, and make no reference to material points that are not contained on this line. 
The general model is geometrically exact and allows an arbitrary number of auxiliary ``warping amplitude" fields, and specializations are obtained that eliminate the additional fields in a general manner. 

Next a section-level constitutive theory is developed in \cref{sec:gensec} which allows for the consideration of material points that do not lie on the reference line.

\cref{sec:timo} is centered on identifying a minimal set of warping shapes which allow for the exact reconstruction of Saint~Venant's classical continuum solutions through the reduced model of \cref{sec:genwarp}. 
A brief recapitulation of Saint Venant's solution is presented in \cref{sec:sv}. 

\begin{Aside}
This draft develops a general family of beam theories enriched by an arbitrary number of warping modes. 
The work begins in the most general setting where we allow for inelasticity and finite deformation. 
We discuss the implications of eliminating warping modes by association with classical strains measures. 
Then we move on to selecting a minimal set of warping shapes. 
To this end, we take a detour into the simple world of the SV class. 
We present Iesan's representation of the class, and break the 3D solution into 5 rigid patterns and 8 warping patterns. 
Next we must determine an association of these patterns with strains arising from the configuration manifold we investigated earlier; ie, 
$$\left\{\bm{u}: [0,L]\rightarrow \mathbb{R}^3, \FrameRotation: [0,L]\rightarrow \mathrm{SO}(3), \bm{\alpha}: [0,L]\rightarrow \mathbb{R}^{n_w}\right\}$$ 
At this point we are confronted with the fundamental ambiguity in associating these fields to full 3D deformation fields. 
So then we introduce the abstraction of traces, $T: \hat{\bm u} \rightarrow (\bm{u}, \bm\theta)$ which formalize this ambiguity. 
There are three natural restrictions we impose on the class of traces we'll investigate: 1) they generate at-most-linear strain measures under the SV class, 2) they preserve plane-section-preserving deformations, and 3) they are capable of exactly reconstructing the exact strain field for any problem in the SV class.
\end{Aside}

\subsection{Details}
We seek a kinematic model for a beam finite element with fully integrated material response, that exactly reproduces the SV solution under linear isotropy. 

Our ultimate goal is to obtain a resultant model that can be used with standard beam finite elements. 
These elements work in terms of the primary equilibrium operator $\mathbb{B}_p$, so if we can return from the section with the quantity that figures in this operator, the element needs no modification. 
It is worth noting that this breaks down under \emph{nonuniform torsion}, because the bishear rate that is picked up by $\mathbf{L}_{\rm vp}\bm{w}'$ is not exactly proportional to any of the 1D strain measures ($\bm\gamma$ and $\bm\kappa$). 
However, this approach will be sufficient to exactly capture the Saint~Venant class which only includes \emph{uniform} torsion.



\subsection{Literature}

% \cite[p39]{vlasov1961thinwalled}

\subsubsection{Shear warping.}

% \paragraph{Elastostatics}

\cite{brown1994nonuniform}


\cite{elfatmi2007nonuniform} put forth a warping beam theory with three warping modes based on the full Saint~Venant solution. 
This work considers  arbitrary section geometry, but neglects the Poisson effect.


Inelastic response with warping is investigated by \cite{neal1961effecta,neal1961effect}.

\subsubsection{Torsion-flexural coupling}
\cite{kim2005exact,elfatmi2007nonuniform}, and \cite{du2020threedimensional}. 
See also \cite{lewinski2021incorporating} who implemented the model of \cite{elfatmi2007nonuniform}.
See also \cite{ladeveze1998new}.


\subsubsection{Element Kinematics}
\paragraph{Finite Deformation}
% Exact rotations plus warping
\cite{simo1991geometricallyexact} introduced a geometrically exact formulation which incorporates one cross-sectional warping mode in a Reissner-like (unconstrained) formulation. 
This was investigated further by \cite{gruttmann2000theory} who emphasized the use of sections intersecting the reference curve at an arbitrary point. 
This theory has also been investigated by \cite{rong2020geometrically} in the context of isogeometric analysis who considered the nonlinear Wagner term. 
However, in this work the reference curve is restricted to intersect sections at the centroid. 

\subsubsection{Section Kinematics}
\paragraph{Moderate Strain (Wagner)}
\cite{rong2020geometrically} incorporated Wagner effects in a geometrically exact beam element with torsion warping.

\cite{chadha2019comprehensive} investigate the a model that accounts for transverse warping induced through axial loading of a beam.

% The asymmetric formulation of \cite{hsiao2000corotational} takes the shear-center as reference point.
